% Tabla A1-T6. Generada automáticamente; no editar a mano.
% Requiere \usepackage{booktabs} (y graphicx si la tabla es ancha).
\begin{table}[htbp]
    \centering
    \caption{Análisis I: reducción del log-loss de cada régimen al añadir la masa estelar, la metalicidad o ambas al modelo base.}
    \label{tab:A1-T6}
    \small
    \resizebox{\textwidth}{!}{%
    \begin{tabular}{lrrrrrrrlrr}
    \toprule
    Régimen & $N_{\rm pl}$ & log-loss B0 & Gan. $M_\star$ & Gan. {[Fe/H]} & Gan. ambas & $M_\star \mid$ {[Fe/H]} & {[Fe/H]} $\mid M_\star$ & Dominante & AUC B0 & AUC B1 \\
    \midrule
    Terrestre & 514 & 1.963 & 0.099 & 0.077 & 0.136 & 0.059 & 0.038 & $M_\star$ & 0.689 & 0.749 \\
    Super-Tierra & 1\,080 & 1.387 & 0.043 & 0.029 & 0.060 & 0.030 & 0.016 & $M_\star$ & 0.568 & 0.646 \\
    Tipo Neptuno & 1\,937 & 0.715 & 0.011 & 0.003 & 0.015 & 0.012 & 0.004 & $M_\star$ & 0.744 & 0.751 \\
    Gigante gaseoso & 720 & 1.857 & 0.088 & 0.152 & 0.192 & 0.039 & 0.103 & {[Fe/H]} & 0.546 & 0.699 \\
    Mayor que Júpiter & 210 & 2.560 & 1.113 & 0.007 & 1.144 & 1.137 & 0.031 & $M_\star$ & 0.812 & 0.945 \\
    \bottomrule
    \end{tabular}
    }
    \par\smallskip
    {\footnotesize\textit{Nota:} Ganancia = log-loss de B0 menos log-loss del modelo ampliado, entre los planetas del régimen. AUC de cada régimen frente al resto (descriptivo).}
\end{table}
